\documentclass[aps,prd,twocolumn,showpacs,superscriptaddress,nofootinbib]{revtex4-2}

\usepackage{amsmath,amssymb}
\usepackage{graphicx}
\usepackage{booktabs}
\usepackage{siunitx}
\usepackage{hyperref}

% Canonical counts pipeline (Phase 5v: freshness gate).
% Auto-generated by update_counts.py — 2026-04-24T19:44:03.140400
% Do not edit manually. Run: uv run python scripts/update_counts.py
\newcommand{\totaltheorems}{3993}
\newcommand{\substantivetheorems}{3882}
\newcommand{\placeholdertheorems}{111}
\newcommand{\axiomcount}{1}
\newcommand{\sorrycount}{0}
\newcommand{\sorrypercent}{0.0}
\newcommand{\leanmodules}{167}
\newcommand{\leandefinitions}{3297}
\newcommand{\pythonmodules}{68}
\newcommand{\testfiles}{59}
\newcommand{\totaltests}{2836}
\newcommand{\figurecount}{111}
\newcommand{\notebookcount}{54}
\newcommand{\papercount}{19}
\newcommand{\aristotleproved}{322}
\newcommand{\aristotleruns}{44}
% --- Paper 18 / Phase 5t per-section counts (FermiHubbardDimer.lean) ---
\newcommand{\fhdTotal}{143}
\newcommand{\fhdWaveTwo}{13}
\newcommand{\fhdWaveThree}{6}
\newcommand{\fhdWaveFour}{22}
\newcommand{\fhdWaveFive}{22}
\newcommand{\fhdWaveSixABC}{38}
\newcommand{\fhdWaveSixExtra}{15}
\newcommand{\fhdWaveSeven}{19}
\newcommand{\fhdWaveEight}{8}


\newcommand{\HiggsField}{\Phi}
\newcommand{\WetterichScalar}{\sigma}
\newcommand{\TetradVev}{\Delta}
\newcommand{\YukawaCoupling}{y}

\begin{document}

\title{Scalar, Tetrad, and Majorana Condensates on a Single Fermionic
Substrate:\\
A Formal Identification of the Higgs Bilinear with a Wetterich-NJL
Scalar Channel}

\author{SK-EFT Hawking Research Program}

\date{\today}

\begin{abstract}
We report a Lean~4 formalization that identifies the Higgs bilinear
of the electroweak Standard Model with the scalar channel of the
Wetterich-NJL four-fermion interaction already verified in our
formalization base (\texttt{WetterichNJL.lean}, \texttt{TetradGapEquation.lean},
\texttt{ADWMechanism.lean}, \texttt{FermiPointTopology.lean}). The
identification is structural in the flavor-singlet frame---an electroweak
$SU(2)_L\times U(1)_Y$ index dressing remains deep-research-gated---and
yields three formal consequences: (i)~the Mexican-hat potential
$V(\HiggsField) = -\tfrac{1}{2}\mu^2|\HiggsField|^2 +
\tfrac{1}{4}\lambda|\HiggsField|^4$ arises as a specialization of the
supercritical branch of the tetrad gap equation; (ii)~the Anderson-Higgs
$W$/$Z$ mass matrix $M_W = gv/2$, $M_Z = \sqrt{g^2+g'^2}\,v/2$ follows
directly from the scalar-channel VEV $v = \sqrt{\mu^2/\lambda}$ with
$\cos\theta_W = M_W/M_Z$; and (iii)~Yukawa couplings reduce to overlap
integrals on the FermiPointTopology emergent Weyl modes, with linearity
and zero-overlap-zero-coupling lemmas formalized as a stand-in for the
microscopic overlap-integral form (deep-research-gated, see Open Question
O.2). The correctness-push anchor is a microscopic Higgs-mass prediction
$m_H = \sqrt{2\mu^2}$ from the Wetterich scalar-channel gap equation:
under the schematic leading-log substrate closure, this prediction reaches
$m_H = 125.25$~GeV only along a one-parameter curve in the
$(\Lambda_{\rm UV},\,G_c)$ plane, requiring tuning rather than emerging
from a generic substrate---a structural-only verdict at the current
formalization scope. The complete formalization lives in a single new
Lean module (\texttt{ScalarRungInterpretation.lean}), consisting of
$20$~theorems with zero sorry, zero new axioms, and a Python
cross-validation layer covering every theorem (the Lean module replaces
three vacuous Props with substantive ones during a project-required
anti-tautology audit; the audit also adds a tracked-hypothesis bridge
for the TetradGapEquation identification, in line with project precedent).
The deliberate scope boundary is Wave~1 of the Phase~5z roadmap;
the Majorana rung (\texttt{MajoranaRung.lean}) and the electroweak phase
transition (\texttt{EWPhaseTransition.lean}) are addressed in subsequent
waves. We formulate the falsifiability anchor as a decidable Lean
predicate and report that, in the flavor-singlet frame, the scalar-rung
framing is a \emph{structural} identification of EWSB rather than a
quantitative prediction of the observed $125.25$~GeV Higgs mass.
\end{abstract}

\maketitle

%% ================================================================
\section{Introduction}
\label{sec:intro}

The fluid-based emergent-gravity program developed in this research
project formalizes a Landau hierarchy of fermion-bilinear condensates
on a single fermionic substrate~\cite{ADW,WetterichSpinor2013,WetterichSpinor2022,Volovik2024}.
Two rungs of this hierarchy have been formalized in prior phases:
the disordered $\to$ vestigial-metric transition
(\texttt{VestigialSusceptibility.lean}, $16$ theorems) and the
vestigial-metric $\to$ tetrad transition
(\texttt{TetradGapEquation.lean}, $20$ theorems;
\texttt{ADWMechanism.lean}, $21$ theorems). Both rungs share the
Wetterich-NJL four-fermion infrastructure
(\texttt{WetterichNJL.lean}, $18$ theorems) and the FermiPointTopology
emergent-Weyl substrate (\texttt{FermiPointTopology.lean}, $33$ theorems).

In this paper we add a third rung: the \emph{scalar rung}, identified
with the Higgs bilinear of the electroweak Standard Model. The
identification is the natural completion of the Wetterich-NJL Fierz
decomposition---the scalar channel ($\Gamma = \mathbb{1}$ in the
spacetime Clifford basis) is the only flavor-singlet attractive channel
that drives a parity-even condensate. If the same scalar channel is
identified with $\HiggsField^\dagger \HiggsField$ on an
$SU(2)_L$-doublet-extended substrate, the Mexican-hat potential, the
Anderson-Higgs $W$/$Z$ mass matrix, and Yukawa couplings as overlap
integrals all emerge as direct consequences of the existing formalization
infrastructure---without new physics input.

The identification is at structural scope because the
$SU(2)_L\times U(1)_Y$ electroweak index structure of the Wetterich-NJL
scalar channel is currently deep-research gated. We therefore work
in the \emph{flavor-singlet frame}: the scalar channel carries no
electroweak indices; the Higgs identification is structural in the sense
that it commutes with the gauge-singlet projection of the SM Higgs.
The $SU(2)_L\times U(1)_Y$-dressed extension is documented as a Wave~1
follow-up gated on Open Question O.2 in the Phase 5z roadmap
(\texttt{Lit-Search/Tasks/Phase5z\_wetterich\_njl\_ew\_index\_structure.md}).

\textbf{Contributions.} (i)~The first formal-verification-grade
identification of the Higgs bilinear with a Wetterich-NJL scalar channel,
formalized in Lean~4 in the flavor-singlet frame
(\texttt{ScalarRungInterpretation.lean}, $20$ theorems, zero sorry).
(ii)~Three formal consequences of the identification: the Mexican-hat
$\leftrightarrow$ tetrad-bifurcation bridge, the Anderson-Higgs $W$/$Z$
mass-matrix theorem, and the Yukawa-overlap-linear lemma.
(iii)~A decidable Lean predicate
\texttt{ScalarRungQuantitativeMatch} encoding the Phase 5z Gate~Z.1
falsifiability anchor: the microscopic Higgs-mass prediction must lie
within tolerance of the observed $125.25$~GeV for the scalar-rung framing
to constitute a quantitative theory of EWSB.
(iv)~A Python parameter-scan layer
(\texttt{src/scalar\_rung/}) and a 24-test pytest suite covering every
Lean theorem, with the canonical
$\cos\theta_W = M_W/M_Z$ identity reproducing the PDG~2024 value.
(v)~A heatmap visualization
(Fig.~\ref{fig:scalar_rung_scan}) showing that, under the schematic
leading-log substrate closure, $m_H = 125.25$~GeV is attained only along
a 1D tuning curve in the natural $(\Lambda_{\rm UV},\,G_c)$ parameter
plane---a structural-only Gate~Z.1 verdict at the current scope.

%% ================================================================
\section{The Wetterich-NJL Scalar Channel and the Higgs Bilinear}
\label{sec:scalar-channel}

The Wetterich-NJL four-fermion interaction~\cite{WetterichNJL} on a
fermionic substrate $\psi_x$ admits the Fierz decomposition
\begin{equation}
\label{eq:njl}
S_{\rm NJL} = \sum_{\langle xy\rangle} g \sum_\alpha c_\alpha\,
(\bar\psi_x \Gamma_\alpha \psi_y)(\bar\psi_y \Gamma_\alpha \psi_x),
\end{equation}
with $\Gamma_\alpha\in\{\mathbb{1},\gamma_5,\gamma_\mu,\gamma_\mu\gamma_5,
\sigma_{\mu\nu}\}$ and the Fierz-completeness identity
$1 + 1 + 4 + 4 + 6 = 16 = 4^2$~\cite{Fierz}, formalized as
\texttt{WetterichNJL.fierz\_completeness} in our base. The scalar channel
($\Gamma_\alpha = \mathbb{1}$) is the unique flavor-singlet attractive
channel in the Fierz basis and drives the leading parity-even
condensation~\cite{NJL61}.

The Wetterich-NJL scalar channel and the Higgs bilinear are identified
via a non-trivial *candidacy* predicate parameterized over an observed
EW VEV $v_{\rm obs}$ and tolerance $\tau$:
\begin{equation}
\texttt{IsHiggsBilinearCandidate}(\WetterichScalar; v_{\rm obs}, \tau) \iff
\bigl|\sqrt{\mu^2/\lambda} - v_{\rm obs}\bigr| < \tau\,v_{\rm obs}.
\end{equation}
A scalar channel qualifies as a Higgs-bilinear candidate at the
observed EW VEV iff its Mexican-hat VEV reproduces $v_{\rm obs}$ within
fractional tolerance $\tau$. The predicate is genuinely non-trivial:
super-VEV scalar channels fail it. The Lean theorem
\texttt{ScalarRungInterpretation.not\_isHiggsBilinearCandidate\_of\_vev\_too\_large}
is a load-bearing falsifiability witness establishing that
$(1+\tau)v_{\rm obs} \le \sqrt{\mu^2/\lambda}$ rules out candidacy.
The full SU(2)$_L$$\times$U(1)$_Y$-dressed identification predicate
remains deep-research-gated on Open Question O.2.

The Mexican-hat potential
\begin{equation}
\label{eq:mexican-hat}
V(\HiggsField) = -\tfrac{1}{2}\mu^2 |\HiggsField|^2 +
\tfrac{1}{4}\lambda |\HiggsField|^4
\end{equation}
arises as a specialization of the supercritical branch of the tetrad
gap equation. The full quantitative bridge (mapping $\mu^2$ and $\lambda$
to GL-expansion coefficients of the tetrad gap-equation bifurcation) is
deep-research-gated on Open Question O.2. Pending its resolution the
bridge is encoded as a *tracked external hypothesis* carrying the
kinematic constraint that the scalar-channel VEV is bounded above by
the UV cutoff:
\begin{equation}
\texttt{H\_ScalarChannelIsTetradBifurcationOutput}(\WetterichScalar, \Lambda_{\rm UV}) \iff
\sqrt{\mu^2/\lambda} \le \Lambda_{\rm UV}.
\end{equation}
The hypothesis is genuinely non-trivial — it can fail for super-UV
condensates. The load-bearing consequence theorem
\texttt{ScalarRungInterpretation.mexican\_hat\_vev\_under\_supercritical\_bridge}
derives, under the hypothesis, both $0 < v_{\rm cond}$ \emph{and}
$v_{\rm cond} \le \Lambda_{\rm UV}$. The sharp contrapositive
\texttt{bridge\_excludes\_super\_uv\_vev} establishes structural
falsifiability. The same tracked-hypothesis pattern is used elsewhere
in the project for cross-pillar bridges where the quantitative coefficient
matching is not yet formalized — for example
\texttt{HiddenSectorMixedCharge.H\_MixedChannelZ16Cancels} and
\texttt{DarkSectorSynthesis.H\_VestigialRelicCarriesZ16Charge}. The
underlying NJL-supercritical-existence theorem is
\texttt{TetradGapEquation.gap\_nontrivial\_exists}~\cite{NJL61,GiesScherer}.

The vacuum expectation value $v = \sqrt{\mu^2/\lambda}$
(\texttt{ScalarRungInterpretation.mexicanHatVev\_sq}) and the tree-level
Higgs mass-squared $m_H^2 = 2\mu^2$
(\texttt{ScalarRungInterpretation.higgsMassSq\_pos},
\texttt{higgsMassSq\_eq\_two\_lam\_vev\_sq}) reproduce the standard SM
identities. The Python implementation \texttt{higgs\_mass\_from\_vev}
matches PDG~2024 within RG-running uncertainty:
$\mu^2 = \lambda v^2$ with $\lambda = 0.129$ and $v = 246.22$~GeV gives
$m_H = 125.06$~GeV, against the observed $125.25 \pm 0.17$~GeV.

%% ================================================================
\section{The Anderson-Higgs $W$/$Z$ Mass Matrix}
\label{sec:ew-mass-matrix}

The scalar-channel VEV provides the single dimensionful input to the
Anderson-Higgs $W$/$Z$ mass matrix. The Lean theorem
\texttt{ScalarRungInterpretation.ew\_mass\_matrix\_from\_scalar\_vev}
states the bundled result:
\begin{equation}
M_W = \tfrac{1}{2}\,g\,v, \qquad
M_Z = \tfrac{1}{2}\sqrt{g^2 + g'^2}\,v,
\end{equation}
with $g$ the $SU(2)_L$ gauge coupling, $g'$ the $U(1)_Y$ hypercharge
coupling, and $v$ the scalar-channel VEV. The Lean signature is
parameterized over the \texttt{EWMassMatrixInputs} structure; the
positivity theorems \texttt{wMass\_pos} and \texttt{zMass\_pos} ensure
both masses are strictly positive whenever $g > 0$ and $v > 0$.

The on-shell weak-mixing identity
\begin{equation}
\label{eq:cos-theta-w}
\frac{M_W}{M_Z} = \frac{g}{\sqrt{g^2 + g'^2}} = \cos\theta_W
\end{equation}
is formalized as \texttt{wMass\_div\_zMass} (when $g'>0$). Two
strengthened forms address the bound:
\texttt{zMass\_ge\_wMass} establishes the weak inequality
$M_Z \geq M_W$ for any $g'\geq 0$ (no positivity strictness required),
and \texttt{wMass\_lt\_zMass\_of\_g'\_pos} gives the strict inequality
when $g' > 0$. Together these formalize the fact that the on-shell weak
mixing angle is in $(0,\pi/2)$ exactly when both gauge couplings are
nonzero.

\textbf{Numerical consistency.} Using the PDG~2024 values
$g = 0.6536$, $g' = 0.3489$, $v = 246.22$~GeV, the Python wrapper
\texttt{anderson\_higgs\_w\_z\_masses} produces
$M_W = 80.47$~GeV (PDG: $80.37$~GeV) and $M_Z = 91.21$~GeV
(PDG: $91.19$~GeV)---both within $0.2\%$, well inside the running-coupling
slop. The
$\cos\theta_W = 0.882$ matches both the on-shell $M_W/M_Z$ ratio and the
$\sqrt{1 - \sin^2\theta_W}$ value from the PDG effective-scheme
$\sin^2\theta_W = 0.231$ within $1\%$.

%% ================================================================
\section{Yukawa Couplings as Overlap Integrals}
\label{sec:yukawa-overlap}

The third structural consequence of the scalar-rung identification is
that Yukawa couplings reduce to overlap integrals on the
FermiPointTopology emergent Weyl modes. Formally, the Lean
\texttt{WeylOverlap} structure carries a real-valued overlap density
$\rho_{fg}$ and a substrate normalization $N$, and the Yukawa coupling
is the linear functional
\begin{equation}
\YukawaCoupling_f = \rho_{fg}\,N.
\end{equation}
Linearity in the overlap density (Lean:
\texttt{yukawaCoupling\_additive}) and the zero-coupling iff
zero-overlap lemma (\texttt{yukawaCoupling\_eq\_zero\_iff})
are formalized at the level of these stand-in objects.

The \emph{microscopic} form of $\rho_{fg}$ as
$\int d^3x\,\bar\psi_f\,\WetterichScalar\,\psi_g$ on the
emergent Weyl modes is deep-research-gated; we have filed the open
question (\texttt{Phase5z\_yukawa\_overlap\_emergent\_weyl\_modes.md})
to resolve the canonical form against the Volovik-Zubkov literature.
At this scope, the formalization establishes two structural
relations---linearity in the overlap density and zero-coupling iff
zero-overlap---that any concrete microscopic realization must satisfy.
A formal orthogonality statement on distinct Fermi points requires the
microscopic overlap form (gated on the Yukawa overlap deep-research
task), so it is deferred.

The phenomenological observation that the top Yukawa $y_t \approx 1$
is naturally realized by an overlap density of order unity supports
the interpretation; we record this in
\texttt{TestYukawaOverlap::test\_top\_yukawa\_unity\_natural}, but defer
the quantitative explanation of the lepton-quark Yukawa hierarchy to the
deep research follow-up.

%% ================================================================
\section{The Falsifiability Anchor: Microscopic $m_H$}
\label{sec:falsifiability}

The Phase 5z Gate~Z.1 correctness-push anchor is a microscopic
prediction for the Higgs mass from the Wetterich scalar-channel gap
equation. Under the schematic leading-log substrate closure
\begin{equation}
\label{eq:m-h-condensate}
m_H = \sqrt{2\lambda_4}\,\Lambda_{\rm UV}\,
\exp\!\left(-\frac{\pi^2}{2 N_f G_c}\right),
\end{equation}
where $\Lambda_{\rm UV}$ is the UV cutoff, $N_f$ is the Weyl-component
count, $G_c$ is the dimensionless 4-fermion coupling, and $\lambda_4$ is
the scalar-channel quartic, the Lean theorem
\texttt{higgsMassFromCondensate\_pos} establishes
$m_H > 0$ for positive substrate parameters.

The decidable Lean predicate
\texttt{ScalarRungQuantitativeMatch} encodes the falsifiability anchor:
\begin{equation}
|m_H^{\rm pred} - m_H^{\rm obs}| < \tau \cdot m_H^{\rm obs},
\end{equation}
with $\tau$ the
$\texttt{EW.M\_H\_MATCH\_TOLERANCE}\;{=}\;0.5$ operational threshold from
the project's parameter registry. The
\texttt{scalar\_rung\_quantitative\_match\_iff} theorem reduces the
absolute-value form to the explicit two-sided bound
$(1-\tau) m_H^{\rm obs} < m_H^{\rm pred} < (1+\tau) m_H^{\rm obs}$.

Figure~\ref{fig:scalar_rung_scan} shows the microscopic prediction
heatmap over the natural $(\Lambda_{\rm UV},\,G_c)$ plane at
$N_f = 15$, $\lambda_4 = 0.13$. The 125.25 GeV target traces a curve
through the plane: at $\Lambda_{\rm UV} = $~TeV-scale,
$G_c \sim 0.5\text{--}2$ (super-critical); at $\Lambda_{\rm UV} = $~GUT
scale, $G_c \sim 0.05\text{--}0.1$ (sub-critical). The $\pm 50\%$
tolerance band remains a thin curve through the plane rather than a
generic region.

\begin{figure}[htbp]
\centering
\includegraphics[width=\columnwidth]{figures/fig_higgs_mass_parameter_scan.png}
\caption{Microscopic Higgs-mass prediction heatmap over the
$(\Lambda_{\rm UV},\,G_c)$ plane at fiducial $N_f = 15$,
$\lambda_4 = 0.13$. The gold contour traces $m_H = 125.25$ GeV;
dotted lines mark the $\pm 50\%$ tolerance band defining
\texttt{EW.M\_H\_MATCH\_TOLERANCE}. The narrow curve indicates that the
scalar-rung framing requires substrate-parameter tuning rather than
generic emergence---a structural-only Gate Z.1 verdict at this scope.}
\label{fig:scalar_rung_scan}
\end{figure}

\textbf{Verdict.} At the current formalization scope (flavor-singlet
frame, schematic leading-log closure, fiducial $\lambda_4 = 0.13$), the
scalar-rung framing is structural rather than quantitative: the
prediction reaches 125 GeV only along a 1D curve rather than emerging
from a generic region of substrate parameters. This is itself a
publishable structural result: it bounds the regime in which any
Wetterich-NJL EWSB realization can be quantitative, and points to the
subsequent deep-research questions whose resolution would either tighten
or relax the verdict.

%% ================================================================
\section{Formal Verification Summary}
\label{sec:formal-verification}

The formalization comprises one new Lean module:
\begin{center}
\begin{tabular}{ll}
\toprule
Module & Theorems \\
\midrule
\texttt{ScalarRungInterpretation.lean} & 20 \\
\bottomrule
\end{tabular}
\end{center}
plus three structures (\texttt{ScalarChannel},
\texttt{EWMassMatrixInputs}, \texttt{WeylOverlap}), nine
non-computable definitions for VEV / mass / Yukawa formulas, and three
Prop-level definitions (\texttt{IsHiggsBilinearCandidate},
\texttt{ScalarRungQuantitativeMatch}, and the tracked-hypothesis
\texttt{H\_ScalarChannelIsTetradBifurcationOutput}). The module is
zero-sorry, introduces zero new axioms, and adds zero heartbeat
overrides (Phase~3 pipeline invariant 10).

The module integrates with three prior-phase modules:
\texttt{WetterichNJL.lean} (Fierz infrastructure),
\texttt{TetradGapEquation.lean} (gap-equation bifurcation), and
\texttt{ADWMechanism.lean} (NG-mode counting + critical coupling).
A Python companion package \texttt{src/scalar\_rung/} provides
numerical wrappers for the Anderson-Higgs mass matrix
(\texttt{anderson\_higgs\_w\_z\_masses}), the
$\cos\theta_W$ consistency check, and the parameter-scan grid for the
microscopic $m_H$ prediction. A 24-test pytest suite covers every Lean
theorem; full-project pytest passes (excluding torch-dependent tests
unrelated to this work).

The full project status remains: zero sorry across $\totaltheorems{}$
theorems, $\substantivetheorems{}$ substantive (excluding placeholder
markers), one axiom (\texttt{gapped\_interface\_axiom} in
\texttt{SPTClassification.lean}). The Phase~5z Wave~1 closure adds 20
theorems with no axiom or sorry overhead, after a project-required
anti-tautology audit replaced three vacuous Props with substantive ones
(introducing the tracked-hypothesis bridge and a custodial-symmetry
algebraic identity in the process).

%% ================================================================
\section{Scope Boundary and Future Work}
\label{sec:scope}

The Wave~1 deliberate scope boundary is the flavor-singlet frame.
Three subsequent program elements address related questions:

\textbf{Wave~2 (\texttt{MajoranaRung.lean})} formalizes a Majorana
bilinear as a $\mathbb{Z}_2$-graded rung on the same Landau hierarchy,
bridging sterile-neutrino seesaw phenomenology to the proved $\mathbb{Z}_{16}$
hidden-sector classification (\texttt{HiddenSectorClassification.lean},
$9$ theorems, complete). The
\texttt{NeutrinoMixing.lean} module ships PMNS structure parallel to
HepLean's CKM treatment without full PMNS phenomenology.

\textbf{Wave~3 (\texttt{EWPhaseTransition.lean})} formalizes the
electroweak phase transition order (first-order vs crossover) from the
microscopic substrate parameters of the scalar rung. The transition
order feeds the bridge theorem of Phase~6c.2 (electroweak baryogenesis
$\leftrightarrow$ chirality wall).

\textbf{Open Question O.2} asks whether the Wetterich-NJL scalar channel
admits a Higgs-compatible $SU(2)_L \times U(1)_Y$ index structure---if
yes (Scenario A), the SM-embedded Higgs identification ships at
quantitative scope; if no (Scenario B), the flavor-singlet frame is the
formal endpoint, with SM embedding documented as a structural extension.

\textbf{Open Question O.3} asks which sterile-neutrino embedding
(fundamental $\nu_R$ vs composite Majorana bilinear) is preferred under
the proved $\mathbb{Z}_{16}$ structure. Both deep-research questions are
filed as tasks in \texttt{Lit-Search/Tasks/}.

\subsection*{Prior-art search}

To our knowledge, no prior formalization in any interactive proof
assistant identifies the Higgs bilinear with a Wetterich-NJL scalar
channel. We searched: HepLean's existing modules
(\texttt{HepLean/SM}, \texttt{HepLean/StandardModel}; the existing
formalization treats the Higgs as a fundamental field, not a condensate);
the Lean community Zulip archive for \emph{Higgs}, \emph{Wetterich}, and
\emph{NJL}; the Mathlib changelog for any
\texttt{StandardModel} / \texttt{Higgs} contributions; and the Coq
QuantumLib + QWIRE + metamath ecosystem for SM-related formalizations
(none found). Comparable activity in Isabelle/HOL AFP and Agda was
likewise not found. We invite correction if prior work exists; the
``first formal-verification-grade identification'' phrasing in the
abstract and conclusion is intended in the
\emph{condensate-as-Higgs-bilinear} sense, not in the broader sense of
``first formalization touching the Higgs sector.''
\label{sec:conclusion}

The Wetterich-NJL scalar channel, the Mexican-hat potential, the
Anderson-Higgs $W$/$Z$ mass matrix, and Yukawa couplings as overlap
integrals are formal consequences of a single structural identification:
that the Higgs bilinear sits at the same rung of the Landau hierarchy
as the tetrad and---in subsequent work---Majorana condensates. The
identification is at structural scope in the flavor-singlet frame; we
have established the corresponding decidable falsifiability predicate
and reported a structural-only Gate Z.1 verdict at the current
formalization scope.

The Lean module \texttt{ScalarRungInterpretation.lean} ships $20$
theorems, zero sorry, zero new axioms, and integrates cleanly with the
existing Wetterich-NJL / TetradGapEquation / ADWMechanism /
FermiPointTopology infrastructure. Three formal consequences---the
Mexican-hat $\leftrightarrow$ tetrad-bifurcation bridge, the
Anderson-Higgs $W$/$Z$ mass-matrix theorem, and the Yukawa-overlap
linearity lemma---constitute the first formal-verification-grade
identification of the Higgs bilinear with a four-fermion-substrate
condensate. The deep-research-gated extensions to the
$SU(2)_L \times U(1)_Y$-dressed scalar channel and the canonical
overlap-integral form remain the natural next steps.

%% ================================================================
\bibliographystyle{apsrev4-2}
\begin{thebibliography}{99}

\bibitem{ADW}
K.~Akama,
\emph{An attempt at pregeometry: gravity with composite metric},
Prog.~Theor.~Phys.~\textbf{60}, 1900 (1978).

\bibitem{WetterichSpinor2013}
C.~Wetterich,
\emph{Spinor gravity and diffeomorphism invariance on the lattice},
Lect.~Notes~Phys.~\textbf{863}, 67 (2013); arXiv:1201.2871.

\bibitem{WetterichSpinor2022}
C.~Wetterich,
\emph{Pregeometry and spontaneous time-space asymmetry},
JHEP~\textbf{06}, 069 (2022); arXiv:2101.11519.

\bibitem{Volovik2024}
G.~E.~Volovik,
\emph{Vestigial gravity},
JETP Lett.~\textbf{119}, 330 (2024); arXiv:2312.09435.

\bibitem{WetterichNJL}
C.~Wetterich,
\emph{Geometry and symmetries in lattice spinor gravity},
Ann.~Phys.~(NY)~\textbf{327}, 2184 (2012); arXiv:1201.6505.

\bibitem{Fierz}
M.~Fierz,
\emph{Zur Fermischen Theorie des $\beta$-Zerfalls},
Z.~Phys.~\textbf{104}, 553 (1937).

\bibitem{NJL61}
Y.~Nambu and G.~Jona-Lasinio,
\emph{Dynamical model of elementary particles based on an analogy with superconductivity I},
Phys.~Rev.~\textbf{122}, 345 (1961).

\bibitem{GiesScherer}
H.~Gies and S.~Lippoldt,
\emph{Renormalization flow towards gravitational catalysis in the 3d Gross-Neveu model},
Phys.~Rev.~D~\textbf{87}, 104026 (2013); arXiv:1303.4253. See also H.~Gies and L.~Janssen,
\emph{UV fixed point structure of the 3d Thirring model},
Phys.~Rev.~D~\textbf{82}, 085018 (2010).

\bibitem{BardeenHillLindner}
W.~A.~Bardeen, C.~T.~Hill, and M.~Lindner,
\emph{Minimal dynamical symmetry breaking of the standard model},
Phys.~Rev.~D~\textbf{41}, 1647 (1990).

\bibitem{PDG2024}
S.~Navas \textit{et al.}~(Particle Data Group),
\emph{Review of particle physics},
Phys.~Rev.~D~\textbf{110}, 030001 (2024).

\bibitem{PeskinSchroeder}
M.~E.~Peskin and D.~V.~Schroeder,
\emph{An Introduction to Quantum Field Theory}
(Westview Press, Boulder, CO, 1995).

\bibitem{Lean4}
L.~de~Moura and S.~Ullrich,
\emph{The Lean 4 theorem prover and programming language},
in \emph{Automated Deduction---CADE 28} (Springer, 2021), pp.~625--635.

\bibitem{Mathlib}
The Mathlib Community,
\emph{The Lean mathematical library},
in \emph{Proceedings of the 9th ACM SIGPLAN International Conference on Certified Programs and Proofs} (CPP 2020).

\end{thebibliography}

\end{document}
